\documentclass[11pt]{article}
\usepackage[margin=1in]{geometry}
\usepackage{amsmath,amssymb,mathtools}
\usepackage{fontspec,unicode-math}
\setmainfont{Latin Modern Roman}
\setsansfont{Latin Modern Sans}
\setmathfont{Latin Modern Math}
\usepackage{xurl,hyperref}
\hypersetup{colorlinks=true,urlcolor=blue,linkcolor=blue}
\newcommand{\sref}[2]{\hyperlink{source:#1:#2}{[S#2]}}
\newcommand{\cataloguescope}[1]{\paragraph{Recorded mathematical question.}#1}
\begin{document}
\section{Percolation criteria for merging planar empires}\label{statement}

\subsection{Definitions and mathematical statement}

\paragraph{Shared notation.}$\mathbb N=\{1,2,\ldots\}$, and $\mathbb Z,\mathbb Q,\mathbb R,\mathbb C$ denote the integers, rationals, reals and complex numbers. Any other conventions are specified locally.


Initially partition $\mathbb R^2$ into polygonal regions, using a regular square, triangular or hexagonal tiling, or a Voronoi tessellation of a planar Poisson process. Call two regions adjacent when they share a boundary segment of positive length. In continuous time, each adjacent pair $A,B$ merges into $A\cup B$ at rate $r(A,B)\ge0$, where the symmetric rate depends only on the geometry of the two regions; the rates are updated as regions merge. Percolation means that an unbounded region exists at some finite time. Nonpercolation means that every region stays bounded at every finite time, almost surely.
\paragraph{Statement recorded in the supplied source.}
Give hypotheses on $r$ sufficient for percolation, and hypotheses sufficient for nonpercolation, specifying which initial tessellations the criteria cover. \sref{statement}{2}
\subsection{Short English statement}
Find geometric merger-rate conditions that cause, or prevent, finite-time growth of an unbounded planar empire.
\subsection{Sources}
\begin{description}
\item[\hypertarget{source:statement:1}{S1}] ULAM.AI and the UnsolvedMath Contributors, \emph{UnsolvedMath: A Curated Collection of Open Mathematics Problems}, catalogue entry. CC BY 4.0. \url{https://huggingface.co/datasets/ulamai/UnsolvedMath}.
\item[\hypertarget{source:statement:2}{S2}] David J. Aldous, \emph{Empires and Percolation}, problem discussed with Vlada Limic in 2001, Process definition and first PROBLEM. \url{https://www.stat.berkeley.edu/~aldous/Research/OP/empires.html}.
\end{description}
\label{end:statement}
\end{document}
